\chapter{Geodesic Geometry and Radome Structure / Geometria Geodésica e Estrutura do Radome}

\section*{English}
The V3 geometry starts from a regular icosahedron with $V_0=12$, $E_0=30$ and $F_0=20$. Each triangular macroface is subdivided into three receiver faces, producing:
\[
V=32,\qquad E=90,\qquad F=60,
\]
with Euler consistency $V-E+F=2$. For preliminary radius $R=\SI{3.1445}{m}$, the closed envelope has diameter $D=2R=\SI{6.2890}{m}$. The radius must not be confused with the height of a closed structure.

The mechanical skeleton, dielectric panels and internal RF modules have different but coordinated geometries. A face is a mechanical and computational sector, not an electromagnetically isolated region. Angular overlap between neighbouring faces is required for calibration, interpolation and direction estimation.

\begin{figure}[H]
\centering
\includegraphics[width=0.95\textwidth]{fig02_zonamento_radome.png}
\caption{Functional zoning / Zoneamento funcional.}
\label{fig:zoning}
\end{figure}

\section*{Português}
A geometria V3 parte de um icosaedro regular com $V_0=12$, $E_0=30$ e $F_0=20$. Cada macroface triangular é subdividida em três faces receptoras, produzindo:
\[
V=32,\qquad E=90,\qquad F=60,
\]
com consistência de Euler $V-E+F=2$. Para raio preliminar $R=\SI{3.1445}{m}$, o envelope fechado possui diâmetro $D=2R=\SI{6.2890}{m}$. O raio não deve ser confundido com a altura de uma estrutura fechada.

O esqueleto mecânico, os painéis dielétricos e os módulos RF internos possuem geometrias diferentes, porém coordenadas. Uma face é um setor mecânico e computacional, não uma região eletromagneticamente isolada. A sobreposição angular entre faces vizinhas é necessária para calibração, interpolação e estimação de direção.
